\chapter{Why You Can Compute with a Qubit}\label{ch:qubit}

\begin{etymology}
\textbf{Abstraction} comes from the Latin \emph{abstrahere}: \emph{ab-},
``away'', and \emph{trahere}, ``to pull''. To abstract is to pull something out
of its context and keep only what matters from one point of view.

\textbf{Superposition} comes from \emph{superponere}: \emph{super-}, ``on top
of'', and \emph{ponere}, ``to put''. Things put on top of each other, occupying
the same place and adding up into one whole.
\end{etymology}

The whole chapter unpacks one sentence:

\begin{quote}
You can compute with qubits because you can encode the candidate answers of a
problem in a superposition, and then apply unitary operations that make the
amplitudes interfere, so that the right answer is favored when you measure.
\end{quote}

Five ideas are in that sentence: the qubit, superposition, interference, unitary
operations and encoding. This chapter takes them one at a time, with as little
mathematics as possible. The next chapters bring the mathematics back.

\section{A qubit is an abstraction}

A qubit is a device with two distinguishable states, which we call $\ket{0}$ and
$\ket{1}$. That is the only thing it shares with a bit. A bit is always one of the
two. A qubit can be in a combination of both, written
\[
  \alpha\ket{0} + \beta\ket{1}.
\]
The numbers $\alpha$ and $\beta$ are \emph{amplitudes}. They are not
probabilities, but they produce them: when you measure, you get $0$ with
probability $|\alpha|^2$ and $1$ with probability $|\beta|^2$, and the two have to
add up to 1.

The usual picture is a coin. On the table, it shows heads or tails. In the air,
you could say it is ``both''. Suppose the coin is weighted so that its state in the
air is $0.866\ket{\text{heads}} + 0.5\ket{\text{tails}}$. Then
$0.866^2 = 0.75$ and $0.5^2 = 0.25$: three times out of four it lands heads. Throw
it a thousand times and the fraction of heads settles near 75\%.

\section{Superposition is not speed}

The coin is a bad picture, and it is worth saying why. A spinning coin is always
either heads or tails at each instant; we just cannot follow it. That is
ignorance, and ignorance is described by ordinary probabilities. A qubit in
superposition is not secretly one of the two. Both components are present at the
same time, and, unlike probabilities, they can cancel each other.

A better image is two waves added together. Take one oscillation and another with
a different frequency, and add them: you get a single wave. Neither of the
originals can be seen in it, but both are in there. This is still an analogy (a
qubit is not a wave in space), so use it only to hold the idea of coexistence.

\section{Interference}

Here is the mechanism, in three steps, with the operation every textbook uses
first. The \emph{Hadamard} operation $H$ acts like this:
\[
  H\ket{0} = \frac{\ket{0} + \ket{1}}{\sqrt{2}}, \qquad
  H\ket{1} = \frac{\ket{0} - \ket{1}}{\sqrt{2}}.
\]
Start from $\ket{0}$ and apply $H$: you get an even superposition, 50/50. Now
flip the sign of the $\ket{1}$ component (a \emph{phase flip}, called $Z$). The
state becomes $(\ket{0} - \ket{1})/\sqrt{2}$. Measure now and you still get 50/50,
because a sign does not change the size of an amplitude. Then apply $H$ again. The
result is $\ket{1}$, with certainty.

Without the phase flip, the same recombination gives $\ket{0}$ with certainty.

\begin{figure}[htbp]
\centering
\begin{tikzpicture}[x=5mm, y=11mm, font=\sffamily\scriptsize]
  \foreach \px/\title/\a/\b in {0/{after $H$}/0.707/0.707, 9/{after the flip}/0.707/-0.707, 18/{after $H$ again}/0/1} {
    \draw[black!40] (\px-0.5,0) -- (\px+5.5,0);
    \fill[hfConcept!70] (\px+0.5,0) rectangle (\px+2,\a);
    \fill[hfHistory!80] (\px+3,0) rectangle (\px+4.5,\b);
    \node[below] at (\px+1.25,{min(0,\a)}) {$\ket{0}$};
    \node[below] at (\px+3.75,{min(0,\b)}) {$\ket{1}$};
    \node[above, font=\sffamily\small] at (\px+2.5,1.25) {\title};
  }
  \node[hf/note] at (2.5,-1.25) {50\% / 50\%};
  \node[hf/note] at (11.5,-1.25) {50\% / 50\%};
  \node[hf/note] at (20.5,-1.25) {0\% / 100\%};
\end{tikzpicture}
\caption{Amplitudes, not probabilities. Flipping the sign of $\ket{1}$ leaves the
probabilities at 50/50, but the second $H$ turns that sign into a certain answer.}
\end{figure}

Why does the sign decide the outcome? Because $H$ adds and subtracts. Applied to
$a\ket{0} + b\ket{1}$, it gives
\[
  \frac{a + b}{\sqrt{2}}\ket{0} + \frac{a - b}{\sqrt{2}}\ket{1}.
\]
When $a$ and $b$ have the same sign, they pile up on $\ket{0}$ and cancel on
$\ket{1}$. When the signs are opposite, the reverse happens. The total
probability is conserved; it is only moved around. This is \emph{interference},
and it is the engine of every quantum algorithm: mark the right answer with a sign
the measurement cannot see, then recombine so the sign turns into probability.

\begin{concept}{Interference}
Amplitudes add like signed (in general, complex) numbers before they become
probabilities. Equal signs reinforce; opposite signs cancel.

\textbf{What it fixes:} an algorithm can steer probability toward the right
answer without ever looking at the intermediate state.

\textbf{What it costs:} the phases must survive the whole computation. Any
contact with the environment scrambles them, and the interference is gone. This
is the reason qubits are so hard to build (Part~\ref{part:hardware}).
\end{concept}

The hard question is how to flip the sign of the right answer without knowing it
in advance. That is what an oracle does, and Chapter~\ref{ch:encoding} returns
to it.

\section{Unitary operations}

$H$ and $Z$ are \emph{unitary} operations. For now, two of their properties
matter. They act on all the components of the superposition at once. And they
are reversible: nothing is thrown away, so you can always undo them. They have to
be. An operation that destroyed information along the way would destroy
amplitudes, and there would be nothing left to interfere. The conservation you
saw above, where 50/50 became 0/100 and the total stayed at 1, is a consequence.
Chapter~\ref{ch:state} says exactly what unitary means.

\section{Encoding}

In the coin, the candidate answers were heads and tails, the two states of the
qubit. A real problem needs more room. One qubit holds 2 candidates, two qubits
hold 4, three hold 8, and $n$ qubits hold $2^n$. Encoding means mapping the
candidates of your problem onto those states, so that all of them can be put into
superposition and take part in the interference. Chapter~\ref{ch:encoding} shows
how it is done for a real problem.

\section{How much does that buy?}

Grover's search finds one marked item among $N = 2^n$ in about $\sqrt{N}$ steps.
A classical machine has to check items one by one; give it 24 cores and it checks
24 at a time.

\begin{center}
\begin{tabular}{rrrrr}
\hline
qubits $n$ & candidates $2^n$ & 24 cores, $N/24$ & Grover, $\sqrt{N}$ & ratio \\
\hline
10 & $1.0 \times 10^{3}$ & 43 & 32 & 1.3 \\
20 & $1.0 \times 10^{6}$ & $4.4 \times 10^{4}$ & $1.0 \times 10^{3}$ & 43 \\
30 & $1.1 \times 10^{9}$ & $4.5 \times 10^{7}$ & $3.3 \times 10^{4}$ & $1.4 \times 10^{3}$ \\
40 & $1.1 \times 10^{12}$ & $4.6 \times 10^{10}$ & $1.0 \times 10^{6}$ & $4.4 \times 10^{4}$ \\
50 & $1.1 \times 10^{15}$ & $4.7 \times 10^{13}$ & $3.4 \times 10^{7}$ & $1.4 \times 10^{6}$ \\
\hline
\end{tabular}
\end{center}

Each extra qubit doubles the classical work and multiplies the quantum work by
only $\sqrt{2}$, so the ratio keeps growing.

The table is also misleading, and you should know how. It compares
\emph{steps}, and a step is not the same thing on both sides. A classical core
checks a candidate in nanoseconds. A Grover step on an error-corrected machine is
a long sequence of logical gates, each of which costs many physical operations
and microseconds of time. In 2021, Ryan Babbush and colleagues at Google worked out
that, with realistic error-correction overheads, a quadratic speedup will not beat
a classical machine on the first generations of fault-tolerant hardware. Square-root
speedups are real, but they are the weakest argument for building quantum
computers.
